% 框架总览
\chapter{统一 DSE 框架：从物理约束到线性规划}\label{chap:overview}

第~\ref{chap:background}章展示了晶圆级交换机的三层物理架构、按位置的七组约束与无阻塞潜能语义。
本章给出本文的核心方法论：如何把这些约束统一到一个线性规划中，
并在给定技术选型下求最大端口无阻塞带宽 $B^*$。

先把结论讲清。本文把 $B^*$ 刻画为期待与约束的函数：
$B^* = f(\text{期待},\ \text{约束})$（式~\eqref{eq:intro-bstar}）。
``期待''是性能侧的无阻塞目标；``约束''是物理侧的资源上限。
我们在给定期待与约束下``尽我们所能''求最大的端口无阻塞带宽。
由此，灵敏度是同一句话的另一面：期待降低或约束放宽，$B^*$ 都提高；
这两个旋钮的定量响应是第~\ref{chap:sensitivity}章的主题。

本章按以下顺序展开：先给两层 DSE 架构（外层选型、内层判定）；
再定义性能约束（排列模式 + 负载包络 + 无阻塞潜能）；
然后按物理位置逐一给出七组约束；
最后组装成``固定 $B$ 的可行性 LP''，用二分搜索求 $B^*$，并以 $B^*$ 的语义收尾。

% 两层DSE架构
\section{两层DSE架构：分离"选型"与"判断"}

在进入具体约束之前，先建立框架的整体结构。我们将DSE流程分为两层。

\textbf{外层：技术选型。}
外层决定离散的架构选择——拓扑参数（Dragonfly的$a,p,h$）、互连标准（UCIe等级、SerDes速率）、
裸片参数（面积$A$、TDP功耗$P$、bump pitch）、冷却方案、裸片布局。
外层是一个枚举/搜索问题：在离散的、可能带上非线性经典约束的候选空间中产生设计点。
外层策略可以是暴力枚举、启发式剪枝或基于agent的智能选型——取决于设计阶段和计算预算。

\textbf{内层：多约束耦合判断。}
内层的任务是在外层固定的技术选型下，回答三个问题：
\begin{enumerate}
    \item 拓扑能否承载目标交换带宽？
    \item 物理资源（bump、布线容量、散热能力）是否足够？
    \item 如果不够，哪个约束是瓶颈、改进它能带来多大收益？
\end{enumerate}

两层解耦的关键在于：\textbf{一旦外层固定了物理参数，内层剩下的所有耦合都发生在一个共同的物理量上——链路负载。}
本章的核心任务就是证明：无论哪类约束、无论建模精度，内层的可行性条件都可以写成关于这一负载向量的线性不等式。

\begin{center}
\begin{tabular}{l|l}
\hline
\textbf{外层（技术选型）} & \textbf{内层（多约束耦合判断）} \\
\hline
拓扑：Dragonfly $(a,p,h)$, Mesh, Torus, \dots & \\
互连标准：UCIe等级, SerDes速率 & 变量：链路负载向量 \\
裸片：$A$, $P$, $p_{\text{bump}}$ & \\
冷却：风冷/液冷/浸没/微流道 & 性能 + 几何 + 功耗 \\
裸片布局：（由外层决定，非内层优化目标）& 三族线性不等式联立求解 \\
\hline
\end{tabular}
\end{center}

% 性能约束
\section{性能约束：排列模式、负载包络与无阻塞潜能}\label{sec:ov-performance}

第一类约束回答：给定端口带宽 $B$，网络在最坏流量模式下是否仍具备无阻塞的\textbf{潜能}？

\subsection{输入：排列模式集合 $\mathcal{R}$}

最坏流量模式由 Birkhoff--von Neumann 定理~\cite{birkhoff1946tres}给出：
双随机流量矩阵多面体的顶点恰为排列矩阵，线性负载在最坏情形下必在排列顶点取到。
因此我们预先固定一组排列需求矩阵 $\{\mathbf{D}^{(r)}\}_{r\in\mathcal{R}}$ 作为\textbf{外生输入}——
$\mathbf{D}$ 不是 LP 的决策变量，网络不能自选流量。
（把 $|\mathcal{R}|$ 从 $n!$ 压到可控规模的群论归约只是降低 LP 规模的实现技巧，见附录 A；
它不改变约束形式，也非本文核心贡献。）

\subsection{负载包络}

对每个模式 $r$，分流变量 $\mathbf{f}^{(r)}$ 满足流守恒，得到模式负载 $\mathbf{L}^{(r)}$；
共享包络 $\mathbf{L}$ 取所有模式的最大值：

\begin{equation}\label{eq:route-flow}
    \forall r \in \mathcal{R}:\quad
    \sum_k f_{ij}^{k,(r)} = D_{ij}^{(r)},\qquad
    \mathbf{L}^{(r)} = \mathcal{P}_r \mathbf{f}^{(r)},\qquad
    \mathbf{L} \;\ge\; \mathbf{L}^{(r)}
\end{equation}

其中 $\mathcal{P}_r$ 为模式 $r$ 的路径 incidence。包络意味着按最坏情形配置物理资源：
所有模式同时以包络负载运行，得到保守的资源与温度上界。
为压出最省资源的包络下界，单独求解性能侧时用目标 $\min\sum_e L_e$：

\begin{equation}\label{eq:valiant-lp}
    \min_{\{\mathbf{f}^{(r)}\}, \mathbf{L}} \; \sum_e L_e
    \quad\text{s.t.}\quad
    \sum_k f_{ij}^{k,(r)} = D_{ij}^{(r)},\;
    \mathbf{L}^{(r)} = \mathcal{P}_r \mathbf{f}^{(r)},\;
    \mathbf{L} \ge \mathbf{L}^{(r)} \;\;\forall r \in \mathcal{R}
\end{equation}

该 LP 只涉及拓扑与路由、与物理参数无关，输出``期待''侧的需求下界 $\mathbf{L}$。
$\mathbf{L}$ 是整张 LP 的枢纽：一端接收 $\mathcal{R}$ 个排列模式的路由需求，
另一端经 $\boldsymbol{\ell} = B\,\mathbf{S}_{\text{bw}}^{-1}\mathbf{L}$ 转化为物理代价。
$\mathbf{L}$ 往上走 $\to$ $\boldsymbol{\ell}$ 涨 $\to$ bump、布线、热约束收紧。

\textbf{关键观察}：性能约束的形式始终是关于 $\mathbf{L}$ 的线性等式/不等式；
$\mathbf{D}^{(r)}$ 外生固定，无阻塞语义由包络 $\mathbf{L}\ge\mathbf{L}^{(r)}$ 承担。

% 物理约束按位置
\section{物理约束按位置}\label{sec:ov-physical}

变量空间与第~\ref{sec:ov-performance}节共享：
$\mathbf{f},\mathbf{L},\boldsymbol{\ell},\mathbf{P},\mathbf{N}^{\text{sig}},\mathbf{N}^{\text{pwr}},\mathbf{T},\mathbf{x}$ 及标量 $B$。
其中 $\boldsymbol{\ell}$ 是唯一的物理桥梁：负载包络经每 lane 带宽换算为物理 lane 数，
lane 数再经每 lane 动态功耗进入功耗：

\begin{equation}\label{eq:lane-def}
    \boldsymbol{\ell} \;=\; B\,\mathbf{S}_{\text{bw}}^{-1}\,\mathbf{L}
    \qquad\Longrightarrow\qquad
    \mathbf{P} \;=\; \mathbf{P}_{peak}(B) + \mathbf{M}\,\mathbf{S}_{\text{dyn}}\,\boldsymbol{\ell}
\end{equation}

其中 $\mathbf{M}$ 为 die-链路 incidence（$M_{v,e}=1 \iff$ 链路 $e$ 的源或宿是 die $v$），
$\mathbf{S}_{\text{bw}},\mathbf{S}_{\text{dyn}}$ 为每 lane 带宽/动态功耗对角阵。
下面按物理位置给七组约束（对应第~\ref{sec:bg-constraints}节的物理描述）。

\subsection{die 内（on-die）：零物理代价}\label{sec:ov-ondie}

\begin{equation}\label{eq:ondie}
    S_{\text{bw},e} = \infty,\qquad S_{\text{dyn},e} = 0 \qquad \forall e \in \mathcal{E}_{\text{on-die}}
\end{equation}

于是 $\ell_e \to 0$，不产生 bump、热、布线需求。

\subsection{die $\leftrightarrow$ interposer 界面：$\mu$bump 的零和竞争}\label{sec:ov-ubump}

每个 die 通过面阵列 $\mu$bump 与 interposer 连接，总焊球数 $N_v^{\text{total}} = \eta\,A_{die}(B)/p^2$。
信号 lane 与电源 lane 共享同一批焊球：

\begin{equation}\label{eq:ubump}
    \mathbf{N}^{\text{sig}} = \mathbf{M}\,\boldsymbol{\ell},\qquad
    \mathbf{N}^{\text{pwr}} = \mathbf{S}_{\text{in}}^{-1}\,\mathbf{P},\qquad
    \mathbf{N}^{\text{sig}} + \mathbf{N}^{\text{pwr}} \;\le\; \mathbf{N}^{\text{total}}(B)
\end{equation}

其中 $\mathbf{S}_{\text{in}}$ 为功率-bump 换算对角阵：$[\mathbf{S}_{\text{in}}]_{vv} = V_{dd}\,I_{bump}$。
物理含义：die 功耗越高，电源 bump 需求越大，留给信号的 bump 越少——性能追求反噬实现可行性。

\subsection{interposer 布线层}\label{sec:ov-wiring}

布线分配变量 $\mathbf{x}$（链路 $\times$ 候选路径）满足：

\begin{equation}\label{eq:wiring}
    \mathbf{A}\,\mathbf{x} = \boldsymbol{\ell},\qquad
    \mathbf{R}\,\mathbf{x} \;\le\; \mathbf{C}
\end{equation}

其中 $\mathbf{A}\mathbf{x}=\boldsymbol{\ell}$ 表达走线需求等于 lane 数；
$\mathbf{R}\mathbf{x}\le\mathbf{C}$ 把边容量、点容量、C4 pad 容量统一为一条容量不等式
（$\mathbf{R}=[\mathbf{R}_{\text{edge}};\ \mathbf{R}_{\text{vert}};\ \mathbf{R}_{\text{pad}}]$）。

\subsection{interposer 热系统}\label{sec:ov-thermal}

稳态热网络把功耗线性映射为温度：

\begin{equation}\label{eq:thermal}
    \mathbf{G}\,\mathbf{T} = \mathbf{P} + \mathbf{b},\qquad
    \mathbf{T} \;\le\; T_{\max}\,\mathbf{1}
\end{equation}

其中 $\mathbf{G}$ 为热导矩阵（对角占优 M-矩阵，$\mathbf{G}^{-1}\ge 0$）。
本文只讨论芯片峰值功耗 $P_{peak}$ 与 PHY 额定动态功耗两个稳态分量；
翘曲、PDN 瞬态、瞬态热、良率、时钟不在本方法讨论范围。

\subsection{interposer $\leftrightarrow$ substrate 界面：C4}\label{sec:ov-c4}

\begin{equation}\label{eq:c4}
    \mathbf{1}^T \boldsymbol{\ell}_{\text{SerDes}} \;\le\; N_{C4}^{\text{SerDes}}
\end{equation}

\subsection{组间（SerDes 全局链路）}\label{sec:ov-intergroup}

组间链路没有独立的约束行——其代价经三个位置体现：
C4（式~\eqref{eq:c4}）、动态功耗（进入式~\eqref{eq:ubump} 与式~\eqref{eq:thermal}）、
pad 容量（进入式~\eqref{eq:wiring}）；拓扑侧需求在性能约束中统一处理。

\subsection{die 缩放模型}\label{sec:ov-diescale}

\begin{equation}\label{eq:die-scale}
    d(B) = d_0 + \alpha_d\,B,\qquad
    A_{die}(B) = d(B)^2,\qquad
    P_{peak}(B) = P_0 + \beta_P\,B
\end{equation}

它使 $\mathbf{N}^{\text{total}}(B)$ 与 $\mathbf{P}_{peak}(B)$ 成为 $B$ 的函数，
从而 bump 预算与热预算随带宽增长而收缩。
（本文数值实验取 $\alpha_d=\beta_P=0$ 的特例；非零时 $\alpha_d,\beta_P$ 正是第~\ref{chap:sensitivity}章物理参数梯度的对象。）

% 统一LP
\section{统一 LP：固定 $B$ 的可行性判定与二分 $B^*$}\label{sec:ov-unified}

现在把所有约束组装起来。给定外层技术选型（所有系数矩阵随之固定）与端口带宽 $B$，
内层判断是否存在一组变量同时满足全部约束。它是纯线性不等式组——\textbf{固定 $B$ 的可行性 LP}：

\begin{equation}\label{eq:unified-lp}
\boxed{
\begin{aligned}
\text{find} \quad & \{\mathbf{f}^{(r)}\}_{r\in\mathcal{R}},\; \{\mathbf{L}^{(r)}\},\; \mathbf{L},\; \boldsymbol{\ell},\; \mathbf{P},\; \mathbf{N}^{\text{sig}},\; \mathbf{N}^{\text{pwr}},\; \mathbf{T},\; \mathbf{x} \\[6pt]
\text{s.t.} \quad & \text{(2.1)}\;\; \textstyle\sum_k f_{ij}^{k,(r)} = D_{ij}^{(r)},\quad
\mathbf{L}^{(r)} = \mathcal{P}_r \mathbf{f}^{(r)},\quad
\mathbf{L} \ge \mathbf{L}^{(r)} \qquad \forall r \\[6pt]
& \text{(2.2)}\;\; S_{\text{bw},e}=\infty,\; S_{\text{dyn},e}=0 \;\; \forall e \in \mathcal{E}_{\text{on-die}} \\[6pt]
& \boldsymbol{\ell} = B\, \mathbf{S}_{\text{bw}}^{-1} \mathbf{L},\qquad
\mathbf{P} = \mathbf{P}_{peak}(B) + \mathbf{M}\mathbf{S}_{\text{dyn}}\boldsymbol{\ell} \\[6pt]
& \text{(2.3)}\;\; \mathbf{N}^{\text{sig}} = \mathbf{M}\boldsymbol{\ell},\quad
\mathbf{N}^{\text{pwr}} = \mathbf{S}_{\text{in}}^{-1}\mathbf{P},\quad
\mathbf{N}^{\text{sig}} + \mathbf{N}^{\text{pwr}} \le \mathbf{N}^{\text{total}}(B) \\[6pt]
& \text{(2.4)}\;\; \mathbf{A}\mathbf{x} = \boldsymbol{\ell},\qquad \mathbf{R}\mathbf{x} \le \mathbf{C} \\[6pt]
& \text{(2.5)}\;\; \mathbf{G}\mathbf{T} = \mathbf{P} + \mathbf{b},\qquad \mathbf{T} \le T_{\max}\mathbf{1} \\[6pt]
& \text{(2.6)}\;\; \mathbf{1}^T \boldsymbol{\ell}_{\text{SerDes}} \le N_{C4}^{\text{SerDes}} \\[6pt]
& \mathbf{f} \ge 0,\;\; \mathbf{L} \ge 0,\;\; \mathbf{x} \ge 0
\end{aligned}
}
\end{equation}

式中 (2.1)--(2.6) 对应性能约束（式~\eqref{eq:route-flow}）与按物理位置的约束
（式~\eqref{eq:ondie}--\eqref{eq:c4}）的编号；die 缩放模型（式~\eqref{eq:die-scale}）
不显式出现，而是经 $\mathbf{N}^{\text{total}}(B)$ 与 $\mathbf{P}_{peak}(B)$ 进入系数与右手边。
组内/组间可分离求解：$B^* = \min(B^*_{\text{intra}}, B^*_{\text{inter}})$——组内用 $\mathcal{E}_{\text{UCIe}}$ + $\mu$bump，
组间用 $\mathcal{E}_{\text{SerDes}}$ + C4。

\subsection{线性化（实现注记）}\label{sec:ov-linearize}

主体约束是物理形式；LP 求解前做一次消元，使全部系数成为常数。
热：消去 $\mathbf{T},\mathbf{P}$（利用 $\mathbf{G}^{-1}\ge 0$）得
$B\,\mathbf{K}\,\mathbf{L}\le\mathbf{rhs}$，其中 $\mathbf{K}=\mathbf{G}^{-1}\mathbf{M}\mathbf{S}_{\text{dyn}}\mathbf{S}_{\text{bw}}^{-1}$，
$\mathbf{rhs}=T_{\max}\mathbf{1}-\mathbf{G}^{-1}(\mathbf{P}_{peak}(B)\mathbf{1}+\mathbf{b})$。
$\mu$bump：动态功耗的 bump 需求折进 lane 系数，静态功耗留在右手边取整。
粗筛（L0）：$\mathbf{1}^T\mathbf{P}\le A_{ip}\,q_{\max}$。
固定 $B$ 后这些系数与右手边都是常数，问题保持线性。

\subsection{二分搜索求 $B^*$}\label{sec:ov-bisection}

可行性关于 $B$ 单调（$B$ 越大物理约束越紧），二分收敛：

\begin{verbatim}
lo 可行, hi 不可行
while hi - lo > step:
    mid = (lo + hi) / 2
    if LP_feasible(mid): lo = mid
    else: hi = mid
return lo
\end{verbatim}

每轮固定 $B$，求解式~\eqref{eq:unified-lp}（无目标函数、纯线性），
$\log_2(B_{\max}/\varepsilon)$ 次 LP 调用即得 $\varepsilon$-近似的 $B^*$。
$B$ 不进入目标函数，意味着每次 LP 规模更小、求解更稳定。

% B* 语义收尾
\section{$B^*$ 的语义收尾}\label{sec:ov-semantics}

回到本章开头的一句话：$B^* = f(\text{期待},\ \text{约束})$。
期待侧（性能）要求 $\mathcal{R}$ 个排列模式下无阻塞，最优自适应路由压出需求下界 $\mathbf{L}$；
约束侧（物理）给出 bump、布线、热、C4 的资源上限（SI 已内嵌于互连标准）。
二分得到的 $B^*$ 刻画``尽我们所能''的最大端口无阻塞带宽——不刻意乐观，也不刻意悲观。

灵敏度随之而来：\textbf{期待降低}（少要求一些流量模式无阻塞，$\mathcal{R}$ 变小）
或\textbf{约束放宽}（工艺、散热、功耗调度改善，rhs 变松），$B^*$ 都提高；
这两个旋钮的定量响应正是第~\ref{chap:sensitivity}章的内容。

